\chapter{Labour Supply and Demand in the Data}\label{ch:b3:labour-supply-and-demand}

Labour is most of what most people sell, and its price is most of what most firms
pay. The Year 1 volume built the supply of hours from the choice between
consumption and leisure, and the demand for labour as a demand derived from the
demand for output. This chapter takes both to the data: who works and how much,
how the answer has changed over a century, how strongly firms respond to the cost
of labour, and how technology changes which tasks people do.

\begin{openingchart}{Women's labour-force participation rate, 1990--2025, six economies
(percent of women aged 15 and over; modelled estimates of the International Labour
Organization).}{World Bank, World Development Indicators (SL.TLF.CACT.FE.ZS)}
\begin{tikzpicture}
\begin{axis}[omtimesten, width=0.72\linewidth, height=0.42\linewidth, xmin=1989,
  xmax=2026, ymin=30, ymax=70, ylabel={Percent}, legend columns=3,
  legend style={font=\footnotesize, at={(0.5,1.03)}, anchor=south}]
  \addplot[omDef, thick, mark=none] table[x=year, y=USA] {figdata/bachelor-3/23-labour-supply-and-demand/out/women.csv};
  \addplot[omThm, thick, mark=none] table[x=year, y=SWE] {figdata/bachelor-3/23-labour-supply-and-demand/out/women.csv};
  \addplot[black, thick, mark=none] table[x=year, y=DEU] {figdata/bachelor-3/23-labour-supply-and-demand/out/women.csv};
  \addplot[omDef, thick, dashed, mark=none] table[x=year, y=JPN] {figdata/bachelor-3/23-labour-supply-and-demand/out/women.csv};
  \addplot[omThm, thick, dashed, mark=none] table[x=year, y=KOR] {figdata/bachelor-3/23-labour-supply-and-demand/out/women.csv};
  \addplot[black, thick, dotted, mark=none] table[x=year, y=ITA] {figdata/bachelor-3/23-labour-supply-and-demand/out/women.csv};
  \legend{United States, Sweden, Germany, Japan, Korea, Italy}
\end{axis}
\end{tikzpicture}
\end{openingchart}

% figdata: women_points
Between 1990 and 2025 the participation of women rose from 45.2 to 55.3~\% in Germany, from
50.1 to 55.9~\% in Japan, from 47.2 to 56.8~\% in Korea and from 35.0 to 40.8~\% in Italy, while it
stayed near 56~\% in the United States and fell slightly in Sweden, from 62.8 to 61.5~\%. The
rate counts all women aged 15 and over, so an ageing population pulls it down even
when women of working age work more; the flat American line hides both movements.
The models of this chapter separate the decision to work from the decision of how
much to work, and ask what moves each.

\begin{recall}
The Year 1 volume derived the supply of hours from the choice between consumption
and leisure, with the income and substitution effects of a wage change and the
labour-supply elasticity; it derived labour demand from production, with factor
shares and the elasticity of substitution. Difference in differences is chapter 3's;
the incidence of a tax wedge is the Year 1 volume's.
\end{recall}

\begin{notation}[Symbols of this chapter]\label{not:b3:labour-supply-and-demand:symbols}
$h$ is hours; $F$ a fixed cost of working; $s$ labour's share of cost; $\sigma$ the elasticity
of substitution between labour and other inputs; $\eta$ the price elasticity of
product demand (positive); $\varepsilon$ a supply elasticity.
\end{notation}

\section{Two margins of labour supply}\label{sec:b3:labour-supply-and-demand:margins}

\begin{definition}[Intensive and extensive margins; Frisch elasticity]\label{def:b3:labour-supply-and-demand:frisch}
\index{intensive margin}\index{extensive margin}\index{Frisch elasticity}
The \emph{intensive margin} of labour supply is the number of hours a worker
supplies; the \emph{extensive margin} is the decision to work at all. The
\emph{Frisch elasticity} is the elasticity of hours with respect to the wage holding
the marginal utility of wealth constant: the response to a wage change expected to
be temporary. The Hicksian (compensated) and Marshallian (uncompensated)
elasticities are the Year 1 volume's.
\end{definition}

\begin{proposition}[A fixed cost makes participation jump]\label{prop:b3:labour-supply-and-demand:fixed-cost}
Let utility be $c - h^2/(2a)$, with $a > 0$, and let a worker earn $wh$ but pay a fixed cost $F > 0$
if $h > 0$. She works $h = aw$ if $w \ge w_R = \sqrt{2F/a}$ and not at all if $w < w_R$; at $w_R$ hours jump from zero to
$\sqrt{2aF}$.
\end{proposition}

\begin{proof}
If she works, she maximises $wh - F - h^2/(2a)$, so $h = aw$ and her net gain is $aw^2/2 - F$. If she does not
work, her gain is zero. She works when $aw^2/2 \ge F$, that is $w \ge \sqrt{2F/a}$. Hours at that wage are $a\sqrt{2F/a}
= \sqrt{2aF} > 0$, and below it zero: no worker chooses a few hours, because the fixed cost
has to be repaid by enough of them.
\end{proof}

Childcare, commuting and the loss of a means-tested benefit are such fixed costs;
\cref{ch:b3:transfers-and-incentives} measured the last. A tax change that raises
take-home pay therefore acts on two margins: those already working may work
longer, and those near the threshold start.

\begin{measure}[Elasticities: what a synthesis finds]
Small frictions --- the cost of noticing a tax change and adjusting hours --- matter
far more on the \omterm{def:b3:labour-supply-and-demand:frisch}{intensive margin} than on the extensive one: the utility cost of
ignoring the tax changes used to estimate hours responses is typically less than
1~\% of earnings. Pooling studies and allowing for such frictions gives Hicksian
elasticities of 0.33 on the \omterm{def:b3:labour-supply-and-demand:frisch}{intensive margin} and 0.25 on the extensive one. % ledger: b3ls:chetty
The frictions can explain why estimates from small, local tax changes are smaller
than those from large, lasting ones.
\end{measure}

\section{Hours over the long run}\label{sec:b3:labour-supply-and-demand:hours}

\begin{omchart}{Average annual hours worked per person employed, 1950--2023, six
economies. Years in which the source's series grows at the same rate year after
year --- interpolated or back-cast, not observed --- are left out: Germany before
1960, Japan before 1953, Korea before 1963 and in 1971--79.}{Penn World Table 11.0
(avh); our screen for constant-growth stretches}
\begin{tikzpicture}
\begin{axis}[omtimesten, width=0.72\linewidth, height=0.42\linewidth, xmin=1948,
  xmax=2025, ymin=1200, ymax=3100, ylabel={Hours a year}, legend columns=3,
  unbounded coords=jump,
  legend style={font=\footnotesize, at={(0.5,1.03)}, anchor=south},
  y tick label style={/pgf/number format/1000 sep={,}}]
  \addplot[omDef, thick, mark=none] table[x=year, y=USA] {figdata/bachelor-3/23-labour-supply-and-demand/out/hours.csv};
  \addplot[black, thick, mark=none] table[x=year, y=DEU] {figdata/bachelor-3/23-labour-supply-and-demand/out/hours.csv};
  \addplot[omThm, thick, mark=none] table[x=year, y=FRA] {figdata/bachelor-3/23-labour-supply-and-demand/out/hours.csv};
  \addplot[omDef, thick, dashed, mark=none] table[x=year, y=GBR] {figdata/bachelor-3/23-labour-supply-and-demand/out/hours.csv};
  \addplot[omThm, thick, dashed, mark=none] table[x=year, y=JPN] {figdata/bachelor-3/23-labour-supply-and-demand/out/hours.csv};
  \addplot[black, thick, dotted, mark=none] table[x=year, y=KOR] {figdata/bachelor-3/23-labour-supply-and-demand/out/hours.csv};
  \legend{United States, Germany, France, United Kingdom, Japan, Korea}
\end{axis}
\end{tikzpicture}
\end{omchart}

% figdata: hours_points
A French worker put in 2{,}351 hours in 1950 and 1{,}487 in 2023; a German one 2{,}181 in 1960,
the first year the source observes, and 1{,}335 in 2023. American hours fell much less,
from 2{,}022 to 1{,}789, and Korean hours rose from 2{,}871 in 1963 to 3{,}040 in 1969 before
falling to 1{,}910. Output per hour worked was multiplied by 10.0 in France and 3.9 in the
United States from 1950, by 7.4 in Germany from 1960 and by 28.5 in Korea from 1963.
% figdata: hours_growth, hours_masked
If hours fall as hourly pay rises, the income effect of the wage has outweighed the
substitution effect: richer workers have bought more leisure. The United States, where
output per hour grew least, cut its hours least --- but Japan, whose output per hour was
multiplied by 11.6 from 1953, cut its hours by a fifth, against two-fifths in France, so
taxes, institutions and preferences matter too.

\section{Women's participation and its causes}\label{sec:b3:labour-supply-and-demand:women}

\begin{omchart}{Ratio of female to male labour-force participation against GDP per
person (PPP, constant dollars, logarithmic scale), 177 countries, 2019, with a
quadratic fit in log income.}{World Bank, World Development Indicators
(SL.TLF.CACT.FM.ZS, NY.GDP.PCAP.PP.KD); our fit}
\begin{tikzpicture}
\begin{semilogxaxis}[omscatter, log ticks with fixed point, width=0.6\linewidth,
  height=0.4\linewidth, xmin=600, xmax=150000, ymin=0, ymax=110,
  xlabel={GDP per person (dollars)}, ylabel={Female/male ratio (\%)}]
  \addplot table[x=y, y=r] {figdata/bachelor-3/23-labour-supply-and-demand/out/ushape.csv};
  \addplot[omThm, thick, no marks, sharp plot] table[x=y, y=r] {figdata/bachelor-3/23-labour-supply-and-demand/out/ushape_fit.csv};
\end{semilogxaxis}
\end{tikzpicture}
\end{omchart}

% figdata: ushape_stats
% ledger: b3ls:goldin
The participation of married women first falls and then rises as countries grow
richer, across countries and in the histories of today's rich ones. As production
moves from the farm and the family business into factories, women's work at home
is no longer counted as work and higher family incomes let them leave the labour
force; once women are educated enough to enter clerical and professional work,
the substitution effect of their own wage dominates. Across countries in 2019 the
fitted U reaches its minimum near 11{,}000 dollars a person, but income explains only
9~\% of the variation ($R^2$ 0.09): the shape is there, and most of the differences
lie elsewhere.

\begin{record}[The pill and women's careers]
The share of American women graduates entering professional programmes ---
law, medicine, business --- rose substantially around 1970, and their age at first
marriage soared just after 1972. The pill, approved in 1960, spread fast among married
women but reached young single women only in the late 1960s, when a series of state
laws lowered the age of majority. The timing, and the variation across states in
when young women gained access, connect the pill to the delay of marriage and to
longer professional training: a lower cost of a long education raised women's
investment in careers. % ledger: b3ls:gk
\end{record}

\section{The elasticity of labour demand}\label{sec:b3:labour-supply-and-demand:demand}

\begin{definition}[Hicks--Marshall rules; scale and substitution effects]\label{def:b3:labour-supply-and-demand:hicks-marshall}
\index{Hicks--Marshall rules}\index{scale effect}\index{substitution effect (labour demand)}
The \emph{Hicks--Marshall rules} state that labour demand is more elastic the easier
it is to substitute other inputs for labour, the more elastic the demand for the
product, the larger labour's share of cost (when product demand is more elastic
than substitution) and the more elastic the supply of other inputs. The
\emph{substitution effect} of a wage rise is the fall in labour at constant output;
the \emph{scale effect} is the fall from the lower output that higher costs bring.
\end{definition}

\begin{proposition}[Hicks--Marshall]\label{prop:b3:labour-supply-and-demand:hm}
With constant returns to scale in labour and one other input, both hired in
competitive markets with the other input's price fixed, and a competitive industry,
the own-wage elasticity of labour demand is $-(1 - s)\sigma - s\eta$.
\end{proposition}

\begin{proof}
Let $K$ be the other input, with price $r$ fixed. Cost minimisation gives $d\ln(L/K) = -\sigma\, d\ln w$. At
constant output, the isoquant gives $s\,d\ln L + (1 - s)\,d\ln K = 0$, with $s$ labour's share of cost;
solving the two equations, $d\ln L = -(1 - s)\sigma\, d\ln w$: the substitution effect. With constant
returns, marginal cost rises by $s\,d\ln w$ and so does the competitive price; output falls
by $\eta s\,d\ln w$, and at a given input ratio labour falls in proportion: the \omterm{def:b3:labour-supply-and-demand:hicks-marshall}{scale effect},
$-s\eta\, d\ln w$. Adding the two gives the result.
\end{proof}

\begin{measure}[How large is the elasticity?]
A \omterm{def:b3:credible-empirical-work:meta-analysis}{meta-analysis} of 942 estimates of the own-wage elasticity of labour demand, from
105 studies, finds that the specification of the model, the data, the sector and the
country explain more than 80~\% of the variation among them, and that publication
favours larger estimates. % ledger: b3ls:lps
There is no single number to quote; there are the \omterm{def:b3:labour-supply-and-demand:hicks-marshall}{Hicks--Marshall rules}, which say
when the number should be large.
\end{measure}

\section{Technology and tasks}\label{sec:b3:labour-supply-and-demand:tasks}

\begin{definition}[Task-based model; routine task; job polarisation]\label{def:b3:labour-supply-and-demand:tasks}
\index{task-based model}\index{routine task}\index{job polarisation}
In a \emph{task-based model}, production uses tasks, each performed by workers of
some skill group or by capital, assigned by \omterm{def:b3:ricardian-trade:comparative-advantage}{comparative advantage}. A \emph{routine
task} follows explicit rules that can be written into a program. \emph{Job
polarisation} is the fall in the employment share of middle-wage occupations
relative to those at the top and the bottom.
\end{definition}

\begin{proposition}[Automating the middle task]\label{prop:b3:labour-supply-and-demand:automation}
Let output be $(y_1y_2y_3)^{1/3}$ over a low, a middle and a high task, each produced one for one by
one unit of a matching group of workers, so that each wage is $1/3$. If capital at price
$\rho$ per unit can perform the middle task and the middle workers move to the low task,
the low and middle groups earn $\sqrt{2/(3\rho)}/6$ and the high group $\sqrt{2/(3\rho)}/3$; with $\rho = 1/4$, wages are
0.272 and 0.544.
\end{proposition}

\begin{proof}
With two workers in the low task and $K$ units of capital, $Y = (2K)^{1/3}$. Each task receives a third
of output, so $\rho K = Y/3$, and $Y^3 = 2Y/(3\rho)$, or $Y = \sqrt{2/(3\rho)}$. The low task's third of output is shared by two
workers: $Y/6$ each; the high group earns $Y/3$. With $\rho = 1/4$, $Y = \sqrt{8/3} = 1.633$, so the low and middle
groups earn 0.272 --- below 1/3 --- and the high group 0.544. The middle workers prefer the
low task's 0.272 to keeping the middle task at capital's price of 0.25.
\end{proof}

\begin{model}{The task-based model of technology}
\assumes Tasks are allocated between capital and groups of workers by \omterm{def:b3:ricardian-trade:comparative-advantage}{comparative
advantage}, and the price of computing power falls steadily.
\predicts Routine middle-wage tasks are automated first; employment moves out of
them toward the top and the bottom; the wages of the top rise.
\tested In the United States, 1980--2005, local labour markets specialised in routine
work adopted information technology more, moved low-skill workers into service
occupations and saw earnings grow at both tails; % ledger: b3ls:ad
in 16 Western European countries, 1993--2010, polarisation was pervasive, and
routine-biased technology together with offshoring explains much of it. % ledger: b3ls:gms
\breaks The wage gains at the bottom need demand for services that cannot be
automated, which the three-task example of the proposition leaves out --- in it the
low wage falls --- and offshoring moves the same tasks abroad for a different reason.
\end{model}

\section{Retirement}\label{sec:b3:labour-supply-and-demand:retirement}

Retirement is the \omterm{def:b3:labour-supply-and-demand:frisch}{extensive margin} at the end of working life. The fixed-cost
proposition applies with a twist: a pension that is withdrawn as earnings rise
raises the implicit tax on working past the pension age, and the \omterm{def:b3:labour-supply-and-demand:frisch}{extensive margin}
responds to it. \Cref{ch:b3:social-insurance} built the pension and its incentives;
in Brenland's survey, no household head aged 65 or over works at all, while 94.9~\% of
those under 65 do (the weekend problem). % figdata: brenland_stats

\section{Exercises}\label{sec:b3:labour-supply-and-demand:exercises}

\begin{exercise}[$\star$]\label{exo:b3:labour-supply-and-demand:1}
In the opening chart, which economy's participation rate rose most between 1990 and
2025? Why can an ageing population lower the rate even when women of every age
work more?
\end{exercise}

\begin{exercise}[$\star$]\label{exo:b3:labour-supply-and-demand:2}
A worker's Hicksian intensive elasticity is 0.33. Her take-home wage rises by 10~\%,
holding utility constant. By how much do her hours change?
\end{exercise}

\begin{exercise}[$\star$]\label{exo:b3:labour-supply-and-demand:3}
Compute the Hicks--Marshall elasticity with $s = 0.3$, $\sigma = 1$ and $\eta = 2$. Split it into the
substitution and \omterm{def:b3:labour-supply-and-demand:hicks-marshall}{scale effects}.
\end{exercise}

\begin{exercise}[$\star$]\label{exo:b3:labour-supply-and-demand:4}
What pattern in a table of occupation shares by decade would contradict the task
model?
\end{exercise}

\begin{exercise}[$\star\star$]\label{exo:b3:labour-supply-and-demand:5}
With the Penn World Table, compute the average annual change in hours per worker
from 1950, or from the first year the second chart shows, to 2023 in its six economies.
\end{exercise}

\begin{exercise}[$\star\star$]\label{exo:b3:labour-supply-and-demand:6}
In the fixed-cost proposition, free childcare lowers $F$. Who enters the labour force,
and what happens to the average hours of those who work?
\end{exercise}

\begin{exercise}[$\star\star$]\label{exo:b3:labour-supply-and-demand:7}
With $a = 100$ and $F = 3{,}200$ dalens a year, find the reservation wage and
the hours worked at it.
\end{exercise}

\begin{exercise}[$\star\star$]\label{exo:b3:labour-supply-and-demand:8}
Fit the U shape of the third chart for 2000 and for 2019. Has the minimum moved?
\end{exercise}

\begin{exercise}[$\star\star$]\label{exo:b3:labour-supply-and-demand:9}
What evidence would contradict a large extensive-margin elasticity among older
workers?
\end{exercise}

\begin{exercise}[$\star\star\star$]\label{exo:b3:labour-supply-and-demand:10}
Derive the Hicks--Marshall formula, and show that a larger labour share makes demand
more elastic if and only if $\eta > \sigma$.
\end{exercise}

\begin{exercise}[$\star\star\star$]\label{exo:b3:labour-supply-and-demand:11}
A payroll-tax cut lowers the wedge between labour cost and take-home pay by 4~\% of
the wage for low-wage workers. Labour demand for them has elasticity 1.0 and supply
0.5. Find the change in take-home pay, in labour cost and in hours.
\end{exercise}

\begin{exercise}[$\star\star\star$]\label{exo:b3:labour-supply-and-demand:12}
In the automation proposition, let $\rho = 1/8$. Compute the wages of the three groups.
Above what price of capital would the middle workers rather keep the middle task at
capital's price than move to the low task?
\end{exercise}

\section{Problem: Brenland's Working Hours}\label{sec:b3:labour-supply-and-demand:problem}

\begin{problem}[{Weekend problem --- Brenland's working hours}]\label{pb:b3:labour-supply-and-demand:hours}
Brenland's household survey (\code{load("households")}) records the hours and gross
income of each household's head; its firm census (\code{load("firms")}) records output
and the wage bill. Brenland employed 4{,}366 thousand people in 2025 (\code{load("accounts")}).
Take the elasticity of substitution between labour and other inputs as 0.7, the
elasticity of product demand as 1.5, and the Hicksian supply elasticities as 0.33 for
hours and 0.25 for participation. % figdata: brenland_stats

\medskip\noindent\textbf{Part I --- hours and wages in the survey.}
\begin{enumerate}
\item What share of heads work no hours? What share of heads under 65?
\item Compute the mean annual hours of those who work.
\item Compute each worker's hourly wage as gross income over hours; find the median.
\item Regress log hours on log wage among workers. Report the slope and its standard
error.
\item Why is this slope not the elasticity of labour supply? Give three reasons.
\item Why does dividing income by hours bias the slope downward?
\end{enumerate}

\medskip\noindent\textbf{Part II --- two margins.} A reform raises every worker's
take-home wage by 4~\%, holding utility constant.
\begin{enumerate}[resume]
\item Predict the change in hours of those already working.
\item Predict the change in the number of workers, in percent and in thousands.
\item Which margin adds more hours?
\item Why might the intensive response in the first year be smaller than predicted?
\item How would you test the prediction with a reform that applied in some regions
only?
\end{enumerate}

\medskip\noindent\textbf{Part III --- labour demand.}
\begin{enumerate}[resume]
\item Compute labour's share of output in the firm census.
\item Using it as the share of cost, compute the Hicks--Marshall elasticity.
\item Split it into substitution and \omterm{def:b3:labour-supply-and-demand:hicks-marshall}{scale effects}.
\item Which effect would be larger in a sector that exports into a competitive world
market?
\item What does the \omterm{def:b3:credible-empirical-work:meta-analysis}{meta-analysis} of the chapter say about using any single value?
\end{enumerate}

\medskip\noindent\textbf{Part IV --- a payroll-tax cut.} Brenland cuts the \omterm{def:b3:efficiency-cost-of-taxation:payroll-tax}{payroll tax}
by 5~\% of the wage. Use the total supply elasticity, intensive plus extensive, and
the Hicksian elasticities as an approximation.
\begin{enumerate}[resume]
\item Compute the rise in take-home pay.
\item Compute the fall in labour cost.
\item Who gains more from the cut, workers or firms, and why?
\item Compute the rise in total hours.
\item Split it into the extensive and intensive parts.
\item Convert the extensive part into thousands of workers.
\item Convert the intensive part into thousands of full-time equivalents.
\item What would change if the cut were temporary?
\item State the result: the employment effect of the payroll-tax cut, intensive plus
extensive, in thousands of workers.
\end{enumerate}
\end{problem}
